\documentclass[10pt,letterpaper]{article}
\usepackage[utf8]{inputenc}
\usepackage[intlimits]{amsmath}
\usepackage{amsfonts}
\usepackage{amssymb}
\usepackage{ragged2e}
\usepackage[letterpaper, margin=1in]{geometry}
\usepackage{graphicx}
\usepackage{cancel}
\usepackage{mathtools}
\usepackage{tabularx}
\usepackage{arydshln}
\usepackage{tensor}
\usepackage{array}
\usepackage{xcolor}
\usepackage[boxed]{algorithm}
\usepackage[noend]{algpseudocode}
\usepackage{listings}
\usepackage{textcomp}
% \usepackage[pdf,tmpdir,singlefile]{graphviz}
\usepackage{mathrsfs}
\usepackage{bbm}
\usepackage{tikz}
\usepackage{tikz-cd}
\usepackage{enumitem}
\usepackage{arydshln}
\usepackage{relsize}
\usepackage{multicol}
\usepackage{scalerel}

\usetikzlibrary{bayesnet}
\setlist{noitemsep}

%%%%%%%%%%%%%%%%%%%%%%%%%%%%%
% Formatting commands
%%%%%%%%%%%%%%%%%%%%%%%%%%%%%
\newcommand{\n}{\hfill\break}
\newcommand{\up}{\vspace{-\baselineskip}}
\newcommand{\hangblock}[2]{\par\noindent\settowidth{\hangindent}{\textbf{#1: }}\textbf{#1: }\!\!\!#2}
\newcommand{\lemma}[1]{\hangblock{Lemma}{#1}}
\newcommand{\defn}[1]{\hangblock{Defn}{#1}}
\newcommand{\thm}[1]{\hangblock{Thm}{#1}}
\newcommand{\cor}[1]{\hangblock{Cor}{#1}}
\newcommand{\prop}[1]{\hangblock{Prop}{#1}}
\newcommand{\ex}[1]{\hangblock{Ex}{#1}}
\newcommand{\exer}[1]{\hangblock{Exer}{#1}}
\newcommand{\fact}[1]{\hangblock{Fact}{#1}}
\newcommand{\remark}[1]{\hangblock{Remark}{#1}}
\newcommand{\proven}{\;$\square$\n}
\newcommand{\problem}[1]{\par\noindent{#1}\n}
\newcommand{\problempart}[2]{\par\noindent\indent{}\settowidth{\hangindent}{\textbf{(#1)} \indent{}}\textbf{(#1) }\!\!\!#2\n}
\newcommand{\ptxt}[1]{\textrm{\textnormal{#1}}}
\newcommand{\inlineeq}[1]{\centerline{$\displaystyle #1$}}
\newcommand{\pageline}{\up\par\noindent\rule{\textwidth}{0.1pt}}

%%%%%%%%%%%%%%%%%%%%%%%%%%%%%
% Math commands
%%%%%%%%%%%%%%%%%%%%%%%%%%%%%
% Set Theory
\newcommand{\card}[1]{\left|#1\right|}
\newcommand{\set}[1]{\left\{#1\right\}}
\newcommand{\setmid}{\;\middle|\;}
\newcommand{\ps}[1]{\mathcal{P}\left(#1\right)}
\newcommand{\pfinite}[1]{\mathcal{P}^{\ptxt{finite}}\left(#1\right)}
\newcommand{\naturals}{\mathbb{N}}
\newcommand{\N}{\naturals}
\newcommand{\integers}{\mathbb{Z}}
\newcommand{\Z}{\integers}
\newcommand{\rationals}{\mathbb{Q}}
\newcommand{\Q}{\rationals}
\newcommand{\reals}{\mathbb{R}}
\newcommand{\R}{\reals}
\newcommand{\complex}{\mathbb{C}}
\newcommand{\C}{\complex}
\newcommand{\halfPlane}{\mathbb{H}}
\let\H\relax
\newcommand{\H}{\halfPlane}
\newcommand{\comp}{^{\complement}}
\DeclareMathOperator{\Hom}{Hom}
\newcommand{\Ind}{\mathbbm{1}}
\newcommand{\cut}{\setminus}
\DeclareMathOperator{\elem}{elem}

% Graph Theory
\let\deg\relax
\DeclareMathOperator{\deg}{deg}
\newcommand{\degp}{\ptxt{deg}^{+}}
\newcommand{\degn}{\ptxt{deg}^{-}}
\newcommand{\precdot}{\mathrel{\ooalign{$\prec$\cr\hidewidth\hbox{$\cdot\mkern0.5mu$}\cr}}}
\newcommand{\succdot}{\mathrel{\ooalign{$\cdot\mkern0.5mu$\cr\hidewidth\hbox{$\succ$}\cr\phantom{$\succ$}}}}
\DeclareMathOperator{\cl}{cl}
\DeclareMathOperator{\affdim}{affdim}

% Probability
\newcommand{\parSymbol}{\P}
\newcommand{\Prob}{\mathbb{P}}
\renewcommand{\P}{\Prob}
\newcommand{\Avg}{\mathbb{E}}
\newcommand{\E}{\Avg}
\DeclareMathOperator{\Var}{Var}
\DeclareMathOperator{\cov}{cov}
\DeclareMathOperator{\Unif}{Unif}
\DeclareMathOperator{\Binom}{Binom}
\newcommand{\CI}{\mathrel{\text{\scalebox{1.07}{$\perp\mkern-10mu\perp$}}}}

% Standard Math
\newcommand{\inv}{^{-1}}
\newcommand{\abs}[1]{\left|#1\right|}
\newcommand{\ceil}[1]{\left\lceil{}#1\right\rceil{}}
\newcommand{\floor}[1]{\left\lfloor{}#1\right\rfloor{}}
\newcommand{\conj}[1]{\overline{#1}}
\newcommand{\of}{\circ}
\newcommand{\tri}{\triangle}
\newcommand{\inj}{\hookrightarrow}
\newcommand{\surj}{\twoheadrightarrow}
\newcommand{\ndiv}{\nmid}
\renewcommand{\epsilon}{\varepsilon}
\newcommand{\divides}{\mid}
\newcommand{\ndivides}{\nmid}
\DeclareMathOperator{\lcm}{lcm}
\DeclareMathOperator{\sgn}{sgn}
\newcommand{\map}[4]{\!\!\!\begin{array}[t]{rcl}#1 & \!\!\!\!\to & \!\!\!\!#2\\ #3 & \!\!\!\!\mapsto & \!\!\!\!#4\end{array}}
\newcommand{\bigsum}[2]{\smashoperator[lr]{\sum_{\scalebox{#1}{$#2$}}}}

% Linear Algebra
\newcommand{\Id}{\textrm{\textnormal{Id}}}
\newcommand{\im}{\textrm{\textnormal{im}}}
\newcommand{\norm}[1]{\abs{\abs{#1}}}
\newcommand{\tpose}{^{T}\!}
\newcommand{\iprod}[1]{\left<#1\right>}
\newcommand{\giprod}{\iprod{\;\,,\;}}
\DeclareMathOperator{\tr}{tr}
\DeclareMathOperator{\trace}{tr}
\newcommand{\chgBasMat}[3]{\!\!\tensor*[_{#1}]{\left[#2\right]}{_{#3}}}
\newcommand{\vecBas}[2]{\tensor*[]{\left[#1\right]}{_{#2}}}
\DeclareMathOperator{\GL}{GL}
\DeclareMathOperator{\Mat}{Mat}
\DeclareMathOperator{\vspan}{span}
\DeclareMathOperator{\rank}{rank}
\newcommand{\V}[1]{\vec{#1}}
\DeclareMathOperator{\proj}{proj}
\DeclareMathOperator{\compProj}{comp}
\DeclareMathOperator{\row}{row}
\newcommand{\smallPMatrix}[1]{\paren{\begin{smallmatrix}#1\end{smallmatrix}}}
\newcommand{\smallBMatrix}[1]{\brack{\begin{smallmatrix}#1\end{smallmatrix}}}
\newcommand{\pmat}[1]{\begin{pmatrix}#1\end{pmatrix}}
\newcommand{\bmat}[1]{\begin{bmatrix}#1\end{bmatrix}}

% Multilinear Algebra
\newcommand{\Lsym}{\L}
\let\L\relax
\DeclareMathOperator{\L}{\mathscr{L}}
\DeclareMathOperator{\A}{\mathcal{A}}
\DeclareMathOperator{\Alt}{Alt}
\DeclareMathOperator{\Sym}{Sym}
\newcommand{\ot}{\otimes}
\newcommand{\ox}{\otimes}
\DeclareMathOperator{\asc}{asc}
\DeclareMathOperator{\asSet}{set}
\DeclareMathOperator{\sort}{sort}
\DeclareMathOperator{\ringA}{\mathring{A}}

% Topology
\newcommand{\closure}[1]{\overline{#1}}
\newcommand{\uball}{\mathcal{U}}
\DeclareMathOperator{\Int}{Int}
\DeclareMathOperator{\Ext}{Ext}
\DeclareMathOperator{\Bd}{Bd}
\DeclareMathOperator{\rInt}{rInt}
\DeclareMathOperator{\ch}{ch}
\DeclareMathOperator{\ah}{ah}
\newcommand{\LargerTau}{\mathlarger{\mathlarger{\mathlarger{\mathlarger{\tau}}}}}
\newcommand{\Tau}{\mathcal{T}}

% Analysis
\DeclareMathOperator{\Graph}{Graph}
\DeclareMathOperator{\epi}{epi}
\DeclareMathOperator{\hypo}{hypo}
\DeclareMathOperator{\supp}{supp}
\newcommand{\lint}[2]{\underset{#1}{\overset{#2}{{\color{black}\underline{{\color{white}\overline{{\color{black}\int}}\color{black}}}}}}}
\newcommand{\uint}[2]{\underset{#1}{\overset{#2}{{\color{white}\underline{{\color{black}\overline{{\color{black}\int}}\color{black}}}}}}}
\newcommand{\alignint}[2]{\underset{#1}{\overset{#2}{{\color{white}\underline{{\color{white}\overline{{\color{black}\int}}\color{black}}}}}}}
\newcommand{\extint}{\ptxt{ext}\int}
\newcommand{\extalignint}[2]{\ptxt{ext}\alignint{#1}{#2}}
\newcommand{\conv}{\ast}
\newcommand{\pd}[2]{\frac{\partial{}#1}{\partial{}#2}}
\newcommand{\del}{\nabla}
\DeclareMathOperator{\grad}{grad}
\DeclareMathOperator{\curl}{curl}
\let\div\relax
\DeclareMathOperator{\div}{div}
\DeclareMathOperator{\vol}{vol}

% Complex Analysis
\let\Re\relax
\DeclareMathOperator{\Re}{Re}
\let\Im\relax
\DeclareMathOperator{\Im}{Im}
\DeclareMathOperator{\Res}{Res}

% Abstract Algebra
\DeclareMathOperator{\ord}{ord}
\newcommand{\generated}[1]{\left<#1\right>}
\newcommand{\cycle}[1]{\smallPMatrix{#1}}
\newcommand{\id}{\ptxt{id}}
\newcommand{\iso}{\cong}
\DeclareMathOperator{\Aut}{Aut}
\DeclareMathOperator{\SL}{SL}
\DeclareMathOperator{\op}{op}
\newcommand{\isom}[4]{\!\!\!\begin{array}[t]{rcl}#1 & \!\!\!\!\overset{\sim}{\to} & \!\!\!\!#2\\ #3 & \!\!\!\!\mapsto & \!\!\!\!#4\end{array}}
\newcommand{\F}{\mathbb{F}}
\newcommand{\acts}{\;\reflectbox{\rotatebox[origin=c]{-90}{$\circlearrowright$}}\;}
\newcommand{\disjunion}{\mathrel{\text{\scalebox{1.07}{$\perp\mkern-10mu\perp$}}}}
\DeclareMathOperator{\SO}{SO}
\DeclareMathOperator{\stab}{stab}
\DeclareMathOperator{\nullity}{nullity}
\DeclareMathOperator{\Perm}{Perm}
\DeclareMathOperator{\nsubgp}{\vartriangleleft}
\DeclareMathOperator{\notnsubgp}{\ntriangleleft}
\newcommand{\presentation}[2]{\left<#1\;\middle|\;#2\right>}

% Convex Optimization
\newcommand{\sectionSymbol}{\S}
\let\S\relax
\newcommand{\S}{\mathbb{S}}
\DeclareMathOperator{\dist}{dist}
\DeclareMathOperator{\dom}{dom}
\DeclareMathOperator{\diag}{diag}
\DeclareMathOperator{\ones}{\mathbbm{1}}
\newcommand{\minimizeOver}[3]{\begin{array}{rl}\underset{#1}{\ptxt{minimize}} & #2\\ \ptxt{subject to} & #3\end{array}}
\newcommand{\maximizeOver}[3]{\begin{array}{rl}\underset{#1}{\ptxt{maximize}} & #2\\ \ptxt{subject to} & #3\end{array}}
\newcommand{\minimizationProblem}[2]{\minimizeOver{}{#1}{#2}}
\newcommand{\maximizationProblem}[2]{\maximizeOver{}{#1}{#2}}
\newcommand{\minimizeOverUnconstrained}[2]{\begin{array}{rl}\underset{#1}{\ptxt{minimize}} & #2\end{array}}
\newcommand{\maximizeOverUnconstrained}[2]{\begin{array}{rl}\underset{#1}{\ptxt{maximize}} & #2\end{array}}
\newcommand{\minimizationUnconstrained}[1]{\minimizeOverUnconstrained{}{#1}}
\newcommand{\maximizationUnconstrained}[1]{\maximizeOverUnconstrained{}{#1}}

% Proofs
\newcommand{\st}{s.t.}
\newcommand{\unique}{!}
\newcommand{\iffdef}{\overset{\ptxt{def}}{\Leftrightarrow}}
\newcommand{\eqdef}{\overset{\ptxt{def}}{=}}
\newcommand{\eqVertical}{\rotatebox[origin=c]{90}{=}}
\newcommand{\mapsfrom}{\mathrel{\reflectbox{\ensuremath{\mapsto}}}}
\newcommand{\mapsdown}{\rotatebox[origin=c]{-90}{$\mapsto$}\mkern2mu}
\newcommand{\mapsup}{\rotatebox[origin=c]{90}{$\mapsto$}\mkern2mu}
\newcommand{\from}{\!\mathrel{\reflectbox{\ensuremath{\to}}}}

% Brackets
\newcommand{\paren}[1]{\left(#1\right)}
\renewcommand{\brack}[1]{\left[#1\right]}
\renewcommand{\brace}[1]{\left\{#1\right\}}
\newcommand{\ang}[1]{\left<#1\right>}

% Algorithms
\algrenewcommand{\algorithmiccomment}[1]{\hskip 1em \texttt{// #1}}
\algrenewcommand\algorithmicrequire{\textbf{Input:}}
\algrenewcommand\algorithmicensure{\textbf{Output:}}
\newcommand{\algP}{\ptxt{\textbf{P}}}
\newcommand{\algNP}{\ptxt{\textbf{NP}}}
\newcommand{\algNPC}{\ptxt{\textbf{NP-Complete}}}
\newcommand{\algNPH}{\ptxt{\textbf{NP-Hard}}}
\newcommand{\algEXP}{\ptxt{\textbf{EXP}}}

%%%%%%%%%%%%%%%%%%%%%%%%%%%%%
% Other commands
%%%%%%%%%%%%%%%%%%%%%%%%%%%%%
\newcommand{\flag}[1]{\textbf{\textcolor{red}{#1}}}
\newcommand{\uSym}{\u}
\let\u\relax
\newcommand{\u}[1]{\underline{#1}}
\newcommand{\bSym}{\b}
\let\b\relax
\newcommand{\b}[1]{\textbf{#1}}
\newcommand{\iSym}{\i}
\let\i\relax
\newcommand{\i}[1]{\textit{#1}}

%%%%%%%%%%%%%%%%%%%%%%%%%%%%%%%%%%%%%%%
% Make l's curvy in math environments %
%%%%%%%%%%%%%%%%%%%%%%%%%%%%%%%%%%%%%%%
\mathcode`l="8000
\begingroup
\makeatletter
\lccode`\~=`\l
\DeclareMathSymbol{\lsb@l}{\mathalpha}{letters}{`l}
\lowercase{\gdef~{\ifnum\the\mathgroup=\m@ne \ell \else \lsb@l \fi}}%
\endgroup

%%%%%%%%%%%%%%%%%%%%%%%%%
% Fix \vdots and \ddots %
%%%%%%%%%%%%%%%%%%%%%%%%%
\usepackage{letltxmacro}
\LetLtxMacro\orgvdots\vdots
\LetLtxMacro\orgddots\ddots

\makeatletter
\DeclareRobustCommand\vdots{%
	\mathpalette\@vdots{}%
}
\newcommand*{\@vdots}[2]{%
	% #1: math style
	% #2: unused
	\sbox0{$#1\cdotp\cdotp\cdotp\m@th$}%
	\sbox2{$#1.\m@th$}%
	\vbox{%
		\dimen@=\wd0 %
		\advance\dimen@ -3\ht2 %
		\kern.5\dimen@
		% remove side bearings
		\dimen@=\wd2 %
		\advance\dimen@ -\ht2 %
		\dimen2=\wd0 %
		\advance\dimen2 -\dimen@
		\vbox to \dimen2{%
			\offinterlineskip
			\copy2 \vfill\copy2 \vfill\copy2 %
		}%
	}%
}
\DeclareRobustCommand\ddots{%
	\mathinner{%
		\mathpalette\@ddots{}%
		\mkern\thinmuskip
	}%
}
\newcommand*{\@ddots}[2]{%
	% #1: math style
	% #2: unused
	\sbox0{$#1\cdotp\cdotp\cdotp\m@th$}%
	\sbox2{$#1.\m@th$}%
	\vbox{%
		\dimen@=\wd0 %
		\advance\dimen@ -3\ht2 %
		\kern.5\dimen@
		% remove side bearings
		\dimen@=\wd2 %
		\advance\dimen@ -\ht2 %
		\dimen2=\wd0 %
		\advance\dimen2 -\dimen@
		\vbox to \dimen2{%
			\offinterlineskip
			\hbox{$#1\mathpunct{.}\m@th$}%
			\vfill
			\hbox{$#1\mathpunct{\kern\wd2}\mathpunct{.}\m@th$}%
			\vfill
			\hbox{$#1\mathpunct{\kern\wd2}\mathpunct{\kern\wd2}\mathpunct{.}\m@th$}%
		}%
	}%
}
\makeatother

\newcommand{\B}{
	\begin{tikzpicture}
	\filldraw [fill=red, draw=black] (0, 0) rectangle (0.37, 0.45);
	\draw [line width=0.5mm, white ] (0.1,0.08) -- (0.1,0.38);
	\draw[line width=0.5mm, white ] (0.1, 0.35) .. controls (0.2, 0.35) and (0.4, 0.2625) .. (0.1, 0.225);
	\draw[line width=0.5mm, white ] (0.1, 0.225) .. controls (0.2, 0.225) and (0.4, 0.1625) .. (0.1, 0.1);
	\end{tikzpicture}
}

\author{Professor Andrew Snowden\\ \small\i{Transcribed by Thomas Cohn}}
\title{Math 493 Lecture 22}
\date{11/20/19} % Can also use \today

\begin{document}
\maketitle
\setlength\RaggedRightParindent{\parindent}
\RaggedRight

\subsection*{Spectral Theorem}

\par\noindent
Let $V$ be a real vector space of dimension $n$ with a positive definite, symmetric bilinear form $\giprod$. Note that positive definite means $\iprod{v,v}>0$ if $v\ne{}0$, and is true iff the bilinear form is non-degenerate and has signature $(n,0)$.\n

\defn{
	A linear operator $T:V\to{}V$ is \textbf{symmetric} if $\iprod{Tv,w}=\iprod{v,Tw}$, $\forall{}v,w\in{}V$.\n
}

\par\noindent
The Spectral Theorem states that a symmetric linear operator is diagonalizable. In particular, all eigenvalues are real.\n

\par\noindent
Let $e_{1},\ldots,e_{n}$ be an orthonormal basis (normal means $\iprod{e_{i},e_{i}}=1$). This exists because of our assumption of the signature. Let $\mathcal{A}$ be the matrix for $T$ on this basis. $V\cong\R^{n}$, $\iprod{v,w}=v\tpose{}\mathcal{A}w$. The symmetric condition on $T$ gives us $\iprod{Tv,w}=\iprod{v,Tw}$, so $(\mathcal{A}v)\tpose{}w=v\tpose(\mathcal{A}w)$, so $v\tpose{}\mathcal{A}\tpose{}w=v\tpose{}\mathcal{A}w$. We conclude $\mathcal{A}=\mathcal{A}\tpose$.\n

\par\noindent
Let $\lambda\in\C$ be an eigenvalue of $\mathcal{A}$, $v\in\C^{n}$ a corresponding eigenvector (so we have $\mathcal{A}v=\lambda{}v$).\n
Write $v=\smallBMatrix{v_{1} & \cdots & v_{n}}\tpose$. Define $\bar v=\smallBMatrix{\bar v_{1} & \cdots & \bar v_{n}}$ (where $\bar \cdot$ denotes complex conjugation). Then
\[
	\bar v\tpose{}v=\bar v_{1}v_{1}+\cdots+\bar v_{n}v_{n}=\abs{v_{1}}^{2}+\cdots+\abs{v_{n}}^{2}\ge{}0
\]
In fact, $\bar v\tpose{}v>0$, because $v\ne{}0$, so $\exists{}i$ \st{} $v_{i}\ne{}0$. Thus,
\[
	\lambda\bar v\tpose{}v=\bar v\tpose\mathcal{A}v=\bar v\tpose\mathcal{A}\tpose{}v=(\mathcal{A}\bar v)\tpose{}v=\conj{(\mathcal{A}v)\tpose}v=\conj{(\lambda{}v)\tpose}v=\conj{\lambda}\bar v\tpose{}v
\]
$\conj{\lambda}\bar v\tpose{}v=\lambda\bar v\tpose{}v$, $\forall{}v$, so because $\bar v\tpose{}v\ne{}0$, $\conj{\lambda}=\lambda$. Thus, $\lambda\in\R$. $T$ has a real eigenvalue $\lambda$, so there is a real eigenvector $v\in\R^{n}$. $V=\vspan(v)\perp\vspan(v)^{\perp}$, so we claim $T$ maps $\vspan(v)^{\perp}$ to itself.\n

\par\noindent
Let $w\in\vspan(v)^{\perp}$, i.e., $\iprod{v,w}=0$. $\iprod{v,Tw}=\iprod{Tv,w}=\iprod{\lambda{}v,w}=\lambda{v,w}=0$. Thus, $Tw\in\vspan(v)^{\perp}$.\n
Let $T'$ be the restrition of $T$ to $\vspan(v)^{\perp}$. $T'$ is symmetric. So by induction on the dimension, spectral theorem applies to $T'$.\n

\par\noindent
Thus, there exists a basis for $\vspan(v)^{\perp}$ consisting of eigenvectors for $T'$. Throwing in $v$, we get a basis for $V$ consisting of eigenvectors of $T$.\proven

\subsection*{Hermitian Forms}

\defn{
	Let $V$ be a $\C$-vector space. A \textbf{Hermitian form} on $V$ is a function $\giprod:V\times{}V\to\C$ such that $\forall{}v,v',w,w'\in{}V,\alpha,\beta\in\C$,
	\begin{enumerate}[leftmargin=4\parindent]
		\item $\giprod$ is $\C$-linear in its second variable: $\iprod{v,\alpha{}w+\beta{}w'}=\alpha\iprod{v,w}+\beta\iprod{v,w'}$.
		\item $\giprod$ is $\C$-sesquilinear in its first variable: $\iprod{\alpha{}v+\beta{}v',w'}=\conj{\alpha}\iprod{v,w}+\conj{\beta}\iprod{v',w}$.
		\item $\iprod{w,v}=\conj{\iprod{v,w}}$. (Note: $\iprod{v,v}\in\R$.)
	\end{enumerate}
}

\defn{
	We say $\giprod$ is \textbf{positive definite} if $\iprod{v,v}>0$ for $v\ne{}0$.\n
}

\ex{
	$V=\C^{n}$, $\iprod{v,w}=\sum_{i=1}^{n}\conj{v_{i}}w_{i}$.\n
	This is a positive definite Hermitian form.\n
}

\subsection*{Anti-Symmetric Forms}

\par\noindent
Let $\giprod$ be a non-degenerate anti-symmetric (that is, $\iprod{v,w}=-\iprod{w,v}$) form on $V$. (And, of course, $V$ is an $F$-vector space, where $F$ is not of characteristic $2$.)\n

\defn{
	A \textbf{symplectic basis} for $V$ is a basis $e_{1},f_{1},\ldots,e_{n},f_{n}$ where
	\begin{itemize}[leftmargin=4\parindent]
		\item $\iprod{e_{i},e_{j}}=0$ if $i\ne{}j$
		\item $\iprod{e_{i},f_{j}}=0$ if $i\ne{}j$
		\item $\iprod{f_{i},f_{j}}=0$ if $i\ne{}j$
		\item $\iprod{e_{i},f_{i}}=1$
	\end{itemize}
	The matrix of $\giprod$ with respect to a symplectic basis is
	\[
		\begin{array}{c}
			\\
			e_{1}\\
			f_{1}\\
			e_{2}\\
			f_{2}\\
			\vdots\\
		\end{array}
		\begin{array}{c}
		\begin{array}{ccccc}
			e_{1} & f_{1} & e_{2} & f_{2} & \cdots
		\end{array}
		\\
		\left[\begin{array}{cc:cc:c}
			0 & -1 & 0 &  0 & \cdots\\
			1 &  0 & 0 &  0 & \cdots\\ \hdashline
			0 &  0 & 0 & -1 & \cdots\\
			0 &  0 & 1 &  0 & \cdots\\ \hdashline
			\vdots & \vdots & \vdots & \vdots & \ddots\\
		\end{array}\right]\\
		\end{array}
	\]
}

\prop{
	$\giprod$ is a non-degenerate, anti-symmetric form on $V$. Then there is a symplectic basis.\n
	Proof: pick $e_{1}\ne{}0$ in $V$. By non-degeneracy, $\exists{}f'_{1}$ \st{} $\iprod{e_{1},f'_{1}}\ne{}0$. So we can scale $f'_{1}$ to $f_{1}$ \st{} $\iprod{e_{1},f_{1}}=1$. Thus, $V=\vspan(e_{1},f_{1})\perp\vspan(e_{1},f_{1})^{\perp}$, and continue by induction.\proven
}

\cor{
	First Corollary: If $\giprod$ is a non-degenerate, anti-symmetric form on $V$, then $\dim{}V$ is even.\n
}

\defn{
	A \textbf{symplectic space} is a pair $(V,\giprod)$, where $\giprod$ is a non-degenerate anti-symmetric form.\n
}

\cor{
	Second Corollary: Any two symplectic spaces of the same dimension are isomorphic.\n
}

\subsection*{Group Representations}

\par\noindent
Let $G$ be a group, $V$ a vector space over a field $F$.\n

\defn{
	A \textbf{linear action} (or \textbf{representation}) of $G$ on $V$ is an action $G\times{}V\to{}V$ such that\break $g\cdot(\alpha{}v+\beta{}w)=\alpha(g\cdot{}v)+\beta(g\cdot{}w)$, $\forall\alpha,\beta\in{}F,v,w\in{}V$.\n
}

\par\noindent
So, given a representation of $G$ on $V$, for each $g\in{}G$, we get an invertible linear map $\map{\rho(g):V}{V}{v}{g\cdot{}v}$\n
This is an element $\rho(g)\in\GL(V)$, so $\map{\rho:G}{\GL(V)}{g}{\rho(g)}$ is a group homomorphism.\n
Conversely, if $\rho:G\to\GL(V)$ is a group homomorphism, then defining $g\cdot{}v=\rho(g)v$ gives a representation of $G$ on $V$.\n

\par\noindent
So, we have a correspondence between representations on $V$ and group homomorphisms $\rho:G\to\GL(V)$.\n

\defn{
	A \textbf{representation} of $G$ is a pair $(V,\rho)$ where $V$ is a vector space and $\rho:G\to\GL(V)$ is a group homomorphism. We often omit $\rho$ from the notation.\n
}

\ex{
	\begin{enumerate}[leftmargin=3\parindent, topsep=0pt]
		\item $G=\GL(V)$, $\rho:G\times\GL(V)\to\GL(V)$ is the identity map.\n
		This is a representation of $G$ on $V$, called the \textbf{standard} representation of $G$.
		\item $G=S_{n}$, $V=F^{n}$ with basis $(e_{1},\ldots,e_{n})$. $\sigma(e_{i})=e_{\sigma(i)}$. $\rho(\sigma)$ is the permutation matrix of $\sigma$.
		\item In general, say a group $G$ acts on a set $X$. Let $F[X]$ be the vector space with basis symbols $[x]$ for $x\in{}X$. An element of $F[X]$ is a formal sum $\sum_{x\in{}X}c_{x}[x]$ where $c_{x}\in{}F$ and only finitely many $c_{x}$'s are nonzero. $F[X]$ is naturally a representation of $G$ by $g\cdot[x]=[gx]$. These are called \textbf{permutation representations}.
		\item $G=D_{n}=\generated{a,b\mid{}a^{2},b^{n},(ab)^{2}}$. Define $\rho:G\to\GL_{2}(\R)$ by $a\mapsto\smallPMatrix{1 & 0\\ 0 & -1}$ and $b\mapsto\smallPMatrix{\cos\frac{2\pi}{n} & -\sin\frac{2\pi}{n}\\ \sin\frac{2\pi}{n} & \cos\frac{2\pi}{n}}$.\n
		This is a two-dimensional representation of $D_{n}$.
		\item $G=\set{1,\sigma}\cong\Z/2\Z$. $V=\C[x,y]$, the set of all polynomials in two variables. $\sigma\cdot{}f(x,y)=f(y,x)$. This is an infinite dimensional representation of $G$.
	\end{enumerate}
}

\end{document}